\documentclass[11pt]{article}
\usepackage[margin=1in]{geometry}
\usepackage[T1]{fontenc}
\usepackage{lmodern}
\usepackage{booktabs}
\usepackage{tabularx}
\usepackage[hidelinks]{hyperref}
\usepackage{url}

\title{Supplementary Protocol: Decision-Time Provenance for Multimodal Fact-Checking}
\author{Sabbir Hossain\\American International University-Bangladesh\\\href{mailto:26-94030-1@student.aiub.edu}{26-94030-1@student.aiub.edu}}
\date{}

\begin{document}
\maketitle

\section{Eligibility and Rating Mapping}
Eligible records are English image--text claims about United States events from 2020--2026. Each record needs an identifiable claim time, an archived or recoverable image, an explicit professional rating, and enough metadata to enforce the temporal cutoff. The mapping accepts three groups of ratings: true or correct; false, fake, or miscaptioned; and unproven, unsubstantiated, or no evidence. Mixed, partially true, satire, opinion, and ambiguous ratings are excluded.

Rating dictionaries are maintained per fact-checking site and frozen before full collection. Evidence sufficiency is a separate audited target and is not inferred solely from the verdict. The 200-case feasibility sample receives a gold sufficiency review under the same decision-time evidence policy; automatically derived targets, if used for training, are marked silver.

\section{Decision-Time Evidence Policy}
The claim time is the earliest documented public appearance. The decision cutoff is the fact-check publication time or a documented equivalent. Each evidence item stores its reported publication time, earliest archive or crawl observation, timestamp source, and confidence level. Confirmed items have consistent metadata and a pre-cutoff independent observation. Observed items have a pre-cutoff observation but incomplete publication metadata. Uncertain items rely on self-reported dates or were first observed after the cutoff.

The primary analysis requires both publication and observation at or before the cutoff and excludes uncertain items. A sensitivity analysis admits uncertain items without changing model parameters. Fact-check pages, fact-check domains, and verdict-bearing post-cutoff material are excluded from the primary evidence pool.

\section{Temporal and Group Splits}
Training uses records from 2020--2023, validation uses 2024, and the final test uses 2025 plus a partial-2026 snapshot ending on 2026-10-24. Results for 2025 and 2026 are reported separately and jointly. The same event, near-duplicate claim, or perceptually identical image cannot cross partitions. Preprocessing parameters and decision thresholds are fitted on training or validation data only.

\section{Experimental Configuration}
\begin{center}
\small
\begin{tabularx}{\textwidth}{p{0.29\textwidth}X}
\toprule
Component & Setting \\
\midrule
Evidence discovery & ClaimReview, GDELT, and Common Crawl metadata \\
Text retrieval & BM25 and BGE-M3 \\
Visual retrieval & SigLIP2 and perceptual hashing \\
Fusion & Reciprocal-rank fusion followed by a BGE reranker \\
Candidate limits & Rerank 100; retain 20 decision-time-valid items \\
Provenance & Typed graph plus bounded dependency-cluster aggregation \\
Decision & Sufficiency gate followed by binary veracity head \\
Calibration & Validation-fitted temperature scaling \\
Abstention & Validation threshold targeting 80\% coverage \\
Repetitions & Five seeds after model selection \\
Selection & 20 deterministic validation-only trials per controlled trainable method \\
Optimization & AdamW, cosine schedule, class-weighted NLL, and early stopping \\
Uncertainty & 10,000 paired event-level bootstrap replicates \\
\bottomrule
\end{tabularx}
\end{center}

\section{Ablations and Stress Tests}
The ablations remove observation-time filtering, dependency clustering, temporal graph features, graph aggregation, textual evidence, visual evidence, the sufficiency gate, and evidence-dropout training. The graph ablation uses parameter-matched cluster pooling. The stress tests inject 1, 5, and 10 syndicated copies, add post-cutoff documents that reveal the verdict, remove the earliest source, introduce dependent and independent conflicts, swap captions, and reuse historical images. An opposite-label donor is used for conflict construction only when pilot review verifies semantic opposition; otherwise a reviewed donor manifest is required. Each intervention stores an exact manifest.

Dependency clustering is validated on manually reviewed evidence pairs and reports precision, recall, and F1. The primary aggregation gives each cluster one unit of base mass, divided among its documents. Publisher reputation is not a primary-model weight because dependence and credibility are different variables.

\section{Dataset Release Policy}
The release contains code, URL and hash metadata, derived graph annotations, retrieval logs, and split manifests. It does not include copyrighted article text or images without explicit permission. Fact-check verdicts remain inherited labels, and automatically extracted evidence links are treated as silver evidence.

\section{Reporting Checklist}
The final submission must report the data snapshot, model and tokenizer revisions, prompts, hardware, precision, search budget, seeds, failures, selection rule, calibration data, decision threshold, elapsed time, and raw prediction locations. It must also report timestamp-confidence counts, sufficiency-label source, dependency-cluster quality, and separate 2025 and partial-2026 results. Every table must be reproducible from preserved result artifacts. At the present stage, the listed software components are implemented, but no real-data result has been measured; synthetic outputs cannot support a scientific claim.

\end{document}

